% ============================================
% 全国大学生数学建模竞赛论文模板
% CUMCM Paper Template
% ============================================
% ============================================
% 国赛 LaTeX 导言区
% 全国大学生数学建模竞赛 (CUMCM)
% ============================================

% === 文档类与中文支持 ===
\documentclass[12pt,a4paper]{article}
\usepackage[UTF8]{ctex}           % 中文支持
\usepackage{fontspec}              % 字体设置

% === 页面布局 ===
\usepackage[top=2.5cm,bottom=2.5cm,left=2.8cm,right=2.8cm]{geometry}
\usepackage{fancyhdr}              % 页眉页脚
\pagestyle{fancy}
\fancyhf{}
\fancyhead[C]{\small 全国大学生数学建模竞赛论文}
\fancyfoot[C]{\thepage}
\renewcommand{\headrulewidth}{0.4pt}

% === 数学宏包 ===
\usepackage{amsmath,amssymb,mathtools}
\usepackage{bm}                    % 粗体数学符号
\usepackage{nicefrac}              % 斜分数

% === 图表 ===
\usepackage{graphicx}              % 插图
\usepackage{float}                 % 浮动体控制
\usepackage{booktabs}              % 三线表
\usepackage{multirow}              % 合并行
\usepackage{longtable}             % 长表格
\usepackage{subcaption}            % 子图

% === 算法 ===
\usepackage[ruled,linesnumbered]{algorithm2e}

% === 列表 ===
\usepackage{enumitem}

% === 颜色与超链接 ===
\usepackage{xcolor}
\usepackage[colorlinks=true,linkcolor=black,citecolor=blue,urlcolor=blue]{hyperref}

% === 参考文献（国标格式）===
\usepackage{gbt7714}
\bibliographystyle{gbt7714-numerical}

% === 代码高亮 ===
\usepackage{listings}
\lstset{
  basicstyle=\small\ttfamily,
  keywordstyle=\color{blue}\bfseries,
  commentstyle=\color{gray},
  stringstyle=\color{red},
  numbers=left,
  numberstyle=\tiny\color{gray},
  frame=single,
  breaklines=true,
  captionpos=b,
  tabsize=4
}

% === 自定义命令 ===
\newcommand{\vect}[1]{\bm{#1}}           % 向量
\newcommand{\mat}[1]{\mathbf{#1}}        % 矩阵
\newcommand{\diff}{\mathrm{d}}           % 微分d
\newcommand{\pd}[2]{\frac{\partial #1}{\partial #2}}  % 偏导数
\newcommand{\dd}[2]{\frac{\diff #1}{\diff #2}}        % 导数

% === 摘要环境 ===
\newenvironment{cnabstract}{
  \begin{center}
    {\large\bfseries 摘\quad 要}
  \end{center}
  \vspace{0.5em}
}{\par\vspace{1em}}

% === 关键词 ===
\newcommand{\cnkeywords}[1]{
  \noindent\textbf{关键词：}#1
}


\begin{document}

% === 标题页 ===
\begin{titlepage}
  \centering
  \vspace*{2cm}
  {\LARGE\bfseries 全国大学生数学建模竞赛论文\par}
  \vspace{1.5cm}
  {\Large\bfseries 论文题目：\underline{\hspace{8cm}}\par}
  \vspace{2cm}

  % --- 摘要 ---
  \begin{cnabstract}
    % 在此填写摘要内容
    % 摘要应包含：问题概述、方法、主要结果
    TODO: 摘要内容
  \end{cnabstract}

  \cnkeywords{关键词1；关键词2；关键词3；关键词4}

  \vfill
\end{titlepage}

% === 目录 ===
\tableofcontents
\newpage

% ========================================
% 一、问题重述与问题分析
% ========================================
\section{问题重述与问题分析}

\subsection{问题重述}
% 用自己的语言重述题目要求，不要照抄原题
% 包含：背景信息、需要解决的具体问题

\subsection{问题分析}
% 对问题进行整体分析
% 包含：问题的数学本质、可能的建模方向、各子问题之间的关系

% ========================================
% 二、模型假设与符号说明
% ========================================
\section{模型假设与符号说明}

\subsection{模型假设}
% 列出建模过程中的合理假设
% 每条假设应有简要说明其合理性
\begin{enumerate}[label=\arabic*.]
  \item 假设1：说明
  \item 假设2：说明
  \item 假设3：说明
\end{enumerate}

\subsection{符号说明}
% 用表格列出主要符号及其含义
\begin{table}[H]
  \centering
  \caption{符号说明}
  \begin{tabular}{cll}
    \toprule
    符号 & 含义 & 单位 \\
    \midrule
    $x$ & 变量说明 & 单位 \\
    $y$ & 变量说明 & 单位 \\
    \bottomrule
  \end{tabular}
\end{table}

% ========================================
% 三、模型建立与公式推导
% ========================================
\section{模型建立与公式推导}

\subsection{模型一：XXX模型}
% 详细描述模型的建立过程
% 包含：数学推导、公式推导、模型的物理/数学意义

% 示例公式
\begin{equation}
  \min_{x} \quad f(x) = \sum_{i=1}^{n} a_i x_i
  \label{eq:model1}
\end{equation}

\subsection{模型二：XXX模型}
% 第二个模型的建立

% ========================================
% 四、模型求解
% ========================================
\section{模型求解}

\subsection{求解算法}
% 描述求解算法的步骤
% 可以用算法环境
\begin{algorithm}[H]
  \caption{算法名称}
  \KwIn{输入参数}
  \KwOut{输出结果}
  初始化\;
  \While{未收敛}{
    步骤1\;
    步骤2\;
  }
  \Return 结果\;
\end{algorithm}

\subsection{求解结果}
% 展示计算结果
% 可以用表格或公式

% ========================================
% 五、结果分析与灵敏度分析
% ========================================
\section{结果分析与灵敏度分析}

\subsection{结果分析}
% 对求解结果进行分析和解释

\subsection{灵敏度分析}
% 分析模型参数变化对结果的影响
% 展示灵敏度分析图表

% ========================================
% 六、专业图表展示与分析
% ========================================
\section{专业图表展示与分析}

% 图表示例
\begin{figure}[H]
  \centering
  % \includegraphics[width=0.8\textwidth]{figures/xxx.png}
  \caption{图表标题}
  \label{fig:example}
\end{figure}

% ========================================
% 七、模型优点、不足与未来应用推广
% ========================================
\section{模型优点、不足与未来应用推广}

\subsection{模型优点}
\begin{enumerate}[label=\arabic*.]
  \item 优点1
  \item 优点2
\end{enumerate}

\subsection{模型不足}
\begin{enumerate}[label=\arabic*.]
  \item 不足1
  \item 不足2
\end{enumerate}

\subsection{未来应用推广}
% 模型的推广方向和应用前景

% ========================================
% 参考文献
% ========================================
\newpage
\begin{thebibliography}{99}
  \bibitem{ref1} 作者. 题目[J]. 期刊, 年份, 卷(期): 页码.
  \bibitem{ref2} 作者. 书名[M]. 出版地: 出版社, 年份.
\end{thebibliography}

% ========================================
% 附录
% ========================================
\newpage
\appendix
\section{附录：代码}

\begin{lstlisting}[language=Python, caption={主要代码}]
# 在此粘贴核心 Python 代码
import numpy as np
import pandas as pd

# 数据处理
# 模型求解
# 结果输出
\end{lstlisting}

\end{document}
